\documentclass[11pt]{article}
\usepackage[a4paper,margin=27mm]{geometry}
\usepackage{amsmath,amssymb,amsthm,mathtools,xcolor,booktabs,array}
\newcommand{\PROVED}{\textcolor{green!40!black}{\textbf{PROVED}}}
\newcommand{\OPEN}{\textcolor{red!70!black}{\textbf{OPEN}}}
\newcommand{\AUDIT}{\textcolor{orange!70!black}{\textbf{AUDIT}}}
\newtheorem{theorem}{Theorem}
\newtheorem{lemma}{Lemma}
\newtheorem{proposition}{Proposition}
\title{Uniform minimizer observability: exact spectral-hole reduction (v130)}
\date{14 August 2026}

\begin{document}
\maketitle

\section{The weight seen by the critical resolvent}

Retain the notation of v124--v129 and put $b=\frac12$.  Define
\begin{align}
 \widehat r_N(\tau)
 &:=\widehat{R_Ne_N}(\tau)
   =\frac{N}{\log N}\frac{\widehat e_N(\tau)}{\tau^2+b^2},       \tag{1}\\
 V_N(\tau)
 &:=\left(\frac{\log N}{N}\right)^2
       (\tau^2+b^2)^2W_N(\tau).                                  \tag{2}
\end{align}
Then the normalization of every minimizing infractor is exactly
\begin{equation}
 \boxed{
 \mathfrak a_N[e_N]
 =\int_{\mathbb R}V_N(\tau)|\widehat r_N(\tau)|^2
       \frac{d\tau}{2\pi}=1.}                                   \tag{3}
\end{equation}
Moreover its Gamma part is
\begin{equation}
 \int_{\mathbb R}g_\Gamma(\tau)V_N(\tau)
       |\widehat r_N(\tau)|^2\frac{d\tau}{2\pi}<1.              \tag{4}
\end{equation}
The all-depth theorem v127 says
\begin{equation}
 V_N(\tau)\longrightarrow C_{\rm W}                              \tag{5}
\end{equation}
locally uniformly on the real line.

\begin{lemma}[Weighted output tightness --- \PROVED]
For every minimizing infractor and every $R$ large enough,
\begin{equation}
 \int_{|\tau|>R}V_N|\widehat r_N|^2\frac{d\tau}{2\pi}
 \le \frac1{\inf_{|\tau|>R}g_\Gamma(\tau)}.                    \tag{6}
\end{equation}
The right side tends to zero independently of $N$.
\end{lemma}

\begin{proof}
This is Markov's inequality applied to (4).  The multiplier
$g_\Gamma$ is even, increasing in $|\tau|$, and tends to infinity.
\end{proof}

Equation (6) is tightness in the weighted norm $L^2(V_Nd\tau)$.
Uniform observability asks for tightness in ordinary $L^2(d\tau)$.

\section{Good frequencies and arithmetic holes}

For $0<\eta<C_{\rm W}/2$ set
\begin{align}
 G_{N,\eta}&:=\{\tau:V_N(\tau)\ge\eta\},                         \tag{7}\\
 H_{N,\eta}&:=\{\tau:V_N(\tau)<\eta\}.                          \tag{8}
\end{align}

\begin{lemma}[The good-frequency contribution --- \PROVED]
Uniformly over all minimizing infractors,
\begin{equation}
 \int_{\{|\tau|>R\}\cap G_{N,\eta}}|\widehat r_N|^2
       \frac{d\tau}{2\pi}
 \le
 \frac1{\eta\inf_{|\tau|>R}g_\Gamma(\tau)}.                   \tag{9}
\end{equation}
\end{lemma}

\begin{proof}
On $G_{N,\eta}$ one has $1\le V_N/\eta$.  Apply (6).
\end{proof}

Thus every possible loss of observability is confined to the explicit
arithmetic sublevel set $H_{N,\eta}$.  Local uniform convergence (5)
also gives, for every fixed $R_0$, an $N_0(R_0,\eta)$ such that
\begin{equation}
 H_{N,\eta}\cap[-R_0,R_0]=\varnothing
 \qquad(N\ge N_0(R_0,\eta)).                              \tag{10}
\end{equation}
The holes, if relevant, must therefore recede to high frequency.

\begin{theorem}[Exact minimizer-hole criterion --- \PROVED]
Uniform minimizer observability is equivalent to the following assertion:
for one (equivalently every) $0<\eta<C_{\rm W}/2$ there is a function
$\varepsilon_\eta(R)\downarrow0$ such that every normalized minimizing
infractor satisfies
\begin{equation}
 \boxed{
 \int_{\{|\tau|>R\}\cap H_{N,\eta}}|\widehat r_N(\tau)|^2
       \frac{d\tau}{2\pi}
 \le\varepsilon_\eta(R).}                               \tag{11}
\end{equation}
\end{theorem}

\begin{proof}
If ordinary $L^2$ tightness holds, (11) follows by restriction.  Conversely,
split $\{|\tau|>R\}$ into $G_{N,\eta}$ and $H_{N,\eta}$.  Lemma 2 controls
the first part by the right side of (9), and (11) controls the second.
Their sum tends to zero independently of $N$.
\end{proof}

This criterion is strictly weaker than asking for a global pointwise
lower bound on $V_N$.  Such a lower bound is not supplied by the existing
estimates: zeros of individual depth polynomials prevent any one-depth
argument, while a simultaneous all-depth lower bound has not been proved.
The criterion asks only that an actual Euler--Lagrange minimizer cannot
place unweighted resolvent mass in the simultaneous all-depth near-zero
set.

\section{What minimization supplies}

Let
\[
 B=\{\tau:g_\Gamma(\tau)<1\}.
\]
For a minimizing infractor with eigenvalue $\lambda_N<1$, the weak
Euler--Lagrange equation of v129 becomes, after the change of variables
(1)--(2),
\begin{equation}
 \int_{\mathbb R}(g_\Gamma-\lambda_N)V_N
       \widehat r_N\overline{\widehat s}
       \frac{d\tau}{2\pi}
 +\widetilde{\mathfrak t}_N(r_N,s)=0                    \tag{12}
\end{equation}
for every $s=R_Nh$ with $h$ in the form domain.  Hence
\begin{align}
 &\int_{B^c}(g_\Gamma-\lambda_N)V_N
       \widehat r_N\overline{\widehat s}\frac{d\tau}{2\pi}
 +\widetilde{\mathfrak t}_N(r_N,s)                       \notag\\
 &\qquad=
 \int_B(\lambda_N-g_\Gamma)V_N
       \widehat r_N\overline{\widehat s}\frac{d\tau}{2\pi}.   \tag{13}
\end{align}
The forcing on the right is confined to the fixed compact band $B$.

\begin{proposition}[The exact estimate still required --- \OPEN]
For every $0<\eta<C_{\rm W}/2$ there should be a function
$\varepsilon_\eta(R)\downarrow0$, independent of $N$, such that every
solution of (12) with $\lambda_N<1$ and (3) satisfies (11).
Equivalently, the compact-band forcing in (13) must have a uniformly
tight response through the positive form on the left, after restriction
to the minimizing eigenspace.
\end{proposition}

The Euler equation does not itself prove this estimate.  On
$H_{N,\eta}$ the coefficient $V_N$ multiplying both sides of (12) is
small.  Dividing by it would use an inverse of the arithmetic leakage
symbol at precisely its near-zero set.  Nor may one infer (11) from the
nonexistence of an infractor: that nonexistence is the conclusion for
which observability is being sought.

\section{A complete fixed-$N$ estimate}

The finite causal inverse from v129 gives, for each fixed $N\ge3$, a
constant $c_N>0$ with
\begin{equation}
 \mathfrak a_N[e]\ge c_N\|e\|_2^2.                              \tag{14}
\end{equation}
Consequently every normalized minimizing infractor satisfies
\begin{equation}
 \boxed{
 \int_{|\tau|>R}|\widehat r_N(\tau)|^2\frac{d\tau}{2\pi}
 \le
 \frac{(N/\log N)^2}{c_N(R^2+\frac14)^2}.}                       \tag{15}
\end{equation}

\begin{proof}
By (1), Plancherel, and (14),
\begin{align*}
 \int_{|\tau|>R}|\widehat r_N|^2\frac{d\tau}{2\pi}
 &\le \frac{(N/\log N)^2}{(R^2+\frac14)^2}\|e_N\|_2^2\\
 &\le \frac{(N/\log N)^2}{c_N(R^2+\frac14)^2}
       \mathfrak a_N[e_N].
\end{align*}
Use $\mathfrak a_N[e_N]=1$.
\end{proof}

This proves observability for every fixed $N$.  It is not uniform in $N$:
the elementary causal inverse provides no bound on
$(N/\log N)^2/c_N$.  Establishing (11) is exactly the improvement that
uses the minimizing equation rather than arbitrary-source deconvolution.

\begin{center}
\begin{tabular}{>{\raggedright\arraybackslash}p{.58\linewidth}
                >{\centering\arraybackslash}p{.14\linewidth}
                >{\raggedright\arraybackslash}p{.18\linewidth}}
\toprule
Statement & Status & Location \\
\midrule
Correct resolvent/output identity & \PROVED & (1)--(4) \\
Weighted output tightness & \PROVED & (6) \\
Unweighted tightness off the arithmetic holes & \PROVED & (9) \\
Exact minimizer-hole equivalence & \PROVED & (11) \\
Observability for each fixed $N$ & \PROVED & (14)--(15) \\
Uniform exclusion of mass inside moving holes & \OPEN & Proposition 4 \\
Uniform minimizer observability, URGT$_1$, $G$, row (d) & \OPEN & after Proposition 4 \\
\bottomrule
\end{tabular}
\end{center}

\end{document}
